\documentclass[10pt]{article}
\usepackage{ctex, amsmath, amssymb, array, booktabs, geometry,enumitem,tasks}
\geometry{a4paper, margin=1.5cm}
\usepackage{titling}\setlength{\droptitle}{-2cm}
\usepackage{tasks}
\usepackage{enumitem}
\setlist[itemize]{itemsep=0pt, parsep=0pt, topsep=3pt}
\setlist[enumerate]{itemsep=0pt, parsep=0pt, topsep=3pt}

\author{学号：\underline{\hspace{4cm}} 姓名：\underline{\hspace{4cm}} }
\title{期权概述与二项式定价公式}
\date{2026年10月10日}
\begin{document}
\maketitle

\begin{enumerate}[label=\textbf{\arabic*.}]

% \section*{5.1.1 期权及有关概念}

\item  {\bf{【计算买入期权与卖出期权的到期价值并判断期权状态】}} %题目 1. 

某股票当前市场价格为 \(S = 60\) 元。某投资者持有：

\begin{enumerate}[label=(\arabic*)]
    \item 一份以该股票为标的资产的\textbf{买入期权}，施权价 \(E = 50\) 元；
    \item 一份以该股票为标的资产的\textbf{卖出期权}，施权价 \(E = 50\) 元。
\end{enumerate}

假设这两份期权现在都已到执行日，请分别计算：

\begin{enumerate}[label=(\arabic*)]
    \item 买入期权在执行日的价值 \(C\)；
    \item 卖出期权在执行日的价值 \(P\)；
    \item 对于买入期权来说，该期权属于币内期权、币上期权还是币外期权？
    \item 对于卖出期权来说，该期权属于币内期权、币上期权还是币外期权？
\end{enumerate}

\vspace{0cm}

\item  {\bf{【根据标的资产价格判断期权状态并计算内在价值】}} %题目 2. 

某投资者关注一份\textbf{卖出期权}，其施权价 \(E = 80\) 元。假设到期日标的资产市场价格 \(S\) 可能出现以下三种情况：

(1) \(S = 95\) 元；
(2) \(S = 80\) 元；
(3) \(S = 65\) 元。

请分别计算这三种情况下卖出期权的执行日价值 \(P\)，并判断该卖出期权分别属于币内期权、币上期权还是币外期权。

\vspace{0cm}

% \section*{5.1.2 期权的一般性质}

\item  {\bf{【利用欧式期权平价公式判断期权价格是否合理】}} %题目 3. 

某不支付红利的股票当前市场价格为 \(S(0) = 39\) 元。有效期为 \(\frac{1}{4}\) 年、施权价为 \(E = 46\) 元的欧式买入期权价格 \(C = 3\) 元，同样施权价和到期日的欧式卖出期权价格 \(P = 2\) 元。年无风险利率为 \(r = 5\%\)。

请利用欧式期权平价公式判断：\(C\) 和 \(P\) 是否符合买入期权、卖出期权平价关系？如果不符合，投资者可以构造怎样的套利组合？当前可得到多少无风险收益？

\vspace{0cm}

\item  {\bf{【计算欧式看涨期权价格的下界并判断是否存在套利】}} %题目 4. 

某不支付红利的股票当前价格 \(S(0) = 50\) 元。欧式看涨期权的施权价 \(E = 45\) 元，到期时间 \(T = 1\) 年，年无风险利率 \(r = 6\%\)。

\begin{enumerate}[label=(\arabic*)]
    \item 根据期权价格边界，计算该欧式看涨期权价格 \(C^E\) 的下界；
    \item 如果市场上该欧式看涨期权价格 \(C^E = 2\) 元，是否低于理论下界？是否存在套利机会？
    \item 如果存在套利机会，请说明基本的套利思路。
\end{enumerate}

\vspace{0cm}

\item  {\bf{【计算欧式看跌期权价格的上界与下界】}} %题目 5. 

某不支付红利的股票当前价格 \(S(0) = 80\) 元。欧式看跌期权的施权价 \(E = 100\) 元，到期时间 \(T = 0.5\) 年，年无风险利率 \(r = 8\%\)。

\begin{enumerate}[label=(\arabic*)]
    \item 计算该欧式看跌期权价格 \(P^E\) 的上界；
    \item 计算该欧式看跌期权价格 \(P^E\) 的下界；
    \item 如果市场上该欧式看跌期权价格 \(P^E = 15\) 元，是否在合理范围内？
\end{enumerate}

\vspace{0cm}

\item  {\bf{【计算期权的时间价值并判断其变化】}} %题目 6. 

某股票当前价格 \(S = 125.23\) 美元。现有两份欧式看涨期权，施权价分别为：

\begin{enumerate}[label=(\arabic*)]
    \item \(E_1 = 110\) 美元，期权价格 \(C_1 = 18.40\) 美元；
    \item \(E_2 = 130\) 美元，期权价格 \(C_2 = 6.78\) 美元。
\end{enumerate}

请分别计算这两份欧式看涨期权的：

(1) 内在价值；\hspace{1cm}
(2) 时间价值；\hspace{1cm}
(3) 判断它们分别属于币内期权、币上期权还是币外期权。

\vspace{0cm}

% \section*{5.1.3 期权的应用举例}

\item  {\bf{【持股购入看跌期权的保值与盈亏平衡点】}} %题目 7. 

某投资者以每股 \(46\) 元的价格买入 100 股股票，同时购入施权价为 \(45\) 元的看跌期权，期权费为 \(1\frac{1}{16}\) 元/份。

\begin{enumerate}[label=(\arabic*)]
    \item 当期权到期日股票价格分别为 \(25\) 元、\(45\) 元、\(65\) 元时，计算该投资者持有股票加看跌期权组合的每股净值；
    \item 计算该策略的盈亏平衡点股票价格。
\end{enumerate}

\vspace{0cm}

\item  {\bf{【空头出售股票并购入看涨期权的保值分析】}} %题目 8. 

某投资者以每股 \(46\) 元的价格空头出售某股票，同时购入施权价为 \(50\) 元的看涨期权，期权费为 \(\frac{7}{16}\) 元。

\begin{enumerate}[label=(\arabic*)]
    \item 当期权到期日股票价格分别为 \(0\) 元、\(25\) 元、\(35\) 元、\(50\) 元、\(65\) 元时，计算该空头出售加购入看涨期权组合的每股净值；
    \item 该策略的最大亏损是多少？出现在什么价格区间？
\end{enumerate}

\vspace{0cm}

\item  {\bf{【出售抵补看涨期权与出售看跌期权的收益分析】}} %题目 9. 

\textbf{情形 A：出售抵补看涨期权}

某投资者以每股 \(46\) 元的价格买入某公司股票 300 股，同时出售 3 份该公司股票的看涨期权，期限 5 个月，施权价为 \(50\) 元，期权费为 \(1\frac{1}{4}\) 元/份。

\begin{enumerate}[label=(\arabic*)]
    \item 计算出售看涨期权获得的期权费总额；
    \item 计算该策略的盈亏平衡点股票价格；
    \item 如果到期日股票价格超过 \(50\) 元，计算投资者每股实际收益。
\end{enumerate}

\textbf{情形 B：出售看跌期权改善买入现状}

某投资者计划以 \(71\) 美元的价格买入股票，投资 \(71\,000\) 美元。他同时出售 3 个月期限的看跌期权，施权价为 \(80\) 美元，期权费每股 \(10.5\) 美元。

\begin{enumerate}[label=(\arabic*),start=4]
    \item 如果期权到期时股票价格低于 \(80\) 美元，计算投资者实际买入股票的总支出和每股成本。
\end{enumerate}

\vspace{0cm}

% \section*{5.1.4 二项式期权定价模型}

\item  {\bf{【推导单期二项式期权定价公式并解释风险中性概率】}} %题目 10. 

考虑单期二项式期权定价模型。设当前股票价格为 \(S\)，在时刻 \(1\) 股票价格有两种可能：

\begin{itemize}
    \item 上涨状态：\(S_u = uS\)，其中 \(u > 1\)；
    \item 下跌状态：\(S_d = dS\)，其中 \(d < 1\)。
\end{itemize}

设无风险利率为 \(r_f\)，令 \(r = 1 + r_f\)。假设市场无套利，因此有：
\(
u > r > d. 
\)

某欧式看涨期权的施权价为 \(E\)，在时刻 \(1\) 的回报分别为：
\(
C_u = \max(uS - E, 0), \quad 
C_d = \max(dS - E, 0). 
\)

请完成以下问题：

\begin{enumerate}[label=(\arabic*)]
    \item 构造一个由 \(\Delta\) 份股票和无风险资产 \(B\) 组成的复制组合，使其在时刻 \(1\) 的价值分别等于 \(C_u\) 和 \(C_d\)，写出方程组；
    \item 解出 \(\Delta\) 和 \(B\) 的表达式；
    \item 利用无套利原理，写出期权在时刻 \(0\) 的价格 \(C\)；
    \item 证明 \(C\) 可以写成如下风险中性定价形式：
    \(\displaystyle 
    C = \frac{1}{r} [p_* C_u + (1 - p_*) C_d],  
    \) 
    并给出 \(p_*\) 的表达式；
    \item 解释为什么 \(p_*\) 被称为风险中性概率，它与股票实际上涨概率 \(q\) 有何区别。
\end{enumerate}

\vspace{0cm}

\item  {\bf{【单期二项式欧式看涨期权定价】}} %题目 11. 

设某股票当前价格 \(S = 20\) 元，单期二项式模型中：
上涨倍数：\(u = 1.2\)；
下跌倍数：\(d = 0.67\)；
无风险利率：\(r_f = 0.1\)，即 \(r = 1.1\)；
欧式看涨期权施权价：\(E = 21\) 元。
请计算：

(1) 时刻 \(1\) 股票上涨和下跌时的价格；\hspace{1cm}
(2) 时刻 \(1\) 期权的回报 \(C_u\) 和 \(C_d\)；

(3) 风险中性概率 \(p_*\)；\hspace{3.4cm}
(4) 期权在时刻 \(0\) 的价格 \(C\)。

\vspace{0cm}

\item  {\bf{【两期二项式欧式看涨期权定价与美式期权提前施权比较】}} %题目 12. 

设某股票当前价格 \(S(0) = 80\) 美元，两期二项式模型中：
上涨倍数：\(u = 1.1\)；
下跌倍数：\(d = 0.95\)；
无风险利率：\(r_f = 0.05\)，即 \(r = 1.05\)；
某美式看跌期权施权价：\(E = 80\) 美元；
到期时间：\(t = 2\)。
请计算：

\begin{enumerate}[label=(\arabic*)]
    \item 风险中性概率 \(p_*\)；
    \item 画出股票价格二叉树；
    \item 计算到期日 \(t = 2\) 时各节点的看跌期权内在价值；
    \item 用连锁法则计算 \(t = 1\) 时各节点美式看跌期权的价值，并判断是否提前施权；
    \item 计算 \(t = 0\) 时美式看跌期权的价格。
\end{enumerate}


\end{enumerate}

\end{document}